\section{Analysis}
\label{sec:analysis}

The aggregate numbers of \cref{tab:main} say how much; this section says where and why. Every stratum is defined from the data alone (\cref{def:strata}) and partitions the test set, so the rows of each table weight to the aggregate.

\subsection{Where the queries are, and where the errors are}

\begin{table}[t]
  \centering\small
  \caption{\textbf{Performance by stratum} on ICEWS18 and ICEWS14s (three seeds, model without competition so that the rows match the ablation base; the competition setting moves the aggregates by 0.25). Strata partition the test set.}
  \label{tab:strata}
  \setlength{\tabcolsep}{2.8pt}\footnotesize
  \begin{tabular}{l rrrr rrrr}
    \toprule
    & \multicolumn{4}{c}{ICEWS18} & \multicolumn{4}{c}{ICEWS14s} \\
    \cmidrule(lr){2-5}\cmidrule(lr){6-9}
    stratum & share & MRR & \hone{} & \hten{} & share & MRR & \hone{} & \hten{} \\
    \midrule
    cold-far   & 17.8\% &  1.84 &  0.64 &  3.18 & 15.9\% &  2.01 &  0.41 &  4.12 \\
    cold-2hop  & 14.0\% &  8.93 &  3.13 & 19.46 & 12.0\% & 11.57 &  3.54 & 28.64 \\
    dyad-only  & 19.2\% & 29.71 & 17.21 & 56.72 & 21.2\% & 37.26 & 23.45 & 66.82 \\
    blocked    & 25.7\% & 34.99 & 16.82 & 72.96 & 16.6\% & 38.43 & 17.64 & 79.93 \\
    clean      & 23.4\% & 85.64 & 76.65 & 99.02 & 34.4\% & 89.93 & 83.32 & 99.34 \\
    \bottomrule
  \end{tabular}
\end{table}

\Cref{tab:strata} is the central table of the paper. On ICEWS18 the clean stratum, where the answer recurred for $(s, r)$ and dominates every distractor on recency and count, is a quarter of the queries and is essentially solved (\hten{} 99.0). The blocked stratum, the same size, is where a monotone copy score provably cannot order the answer first; \model{} ranks it first 16.8\% of the time and in the top ten 73\% of the time, so the ordering information it uses is not count-and-recency. The dyad-only stratum, a fifth of the queries and invisible to every copy mechanism, reaches MRR 29.7 and \hten{} 56.7: the untrained blind-dyad counter reached 20.0 and 42.3 on the same stratum (\cref{sec:diag}), and the learned stream adds ten points on top of it. The two cold strata, a third of the queries, are where the model is weak: an answer that is a partner of a partner is found in the top ten one time in five, and an answer that is further than two hops from the subject is almost never found. The \hten{} gap to the global-graph family in \cref{tab:main} is of the order of the mass in these two rows, and we take it to be where the next step lies.

\subsection{Ablations}

\begin{table}[t]
  \centering\small
  \caption{\textbf{Ablations} on an event dataset and a yearly dataset: MRR and its change relative to the model without competition. Rows marked $^{3}$ are three-seed means (spread at most 0.19); the others are seed 1. The seed spread of the base is 0.04 on ICEWS18 and 0.20 on YAGO. \hone{} moves with MRR in every row (ICEWS18 base 26.10, without the dyadic branch 20.51; YAGO base 90.57, without it 61.37).}
  \label{tab:ablation}
  \setlength{\tabcolsep}{4pt}\footnotesize
  \begin{tabular}{l rr rr}
    \toprule
    & \multicolumn{2}{c}{ICEWS18} & \multicolumn{2}{c}{YAGO} \\
    \cmidrule(lr){2-3}\cmidrule(lr){4-5}
    removed & MRR & $\Delta$ & MRR & $\Delta$ \\
    \midrule
    nothing (base)$^{3}$           & 36.29 &         & 91.82 & \\
    popularity field               & 36.37 & $+0.08$ & 91.81 & $-0.01$ \\
    candidate context              & 36.27 & $-0.02$ & 91.75$^{3}$ & $-0.07$ \\
    path intensity                 & 35.92 & $-0.37$ & 91.70$^{3}$ & $-0.12$ \\
    relation types in the stream   & 35.86 & $-0.43$ & 91.79 & $-0.03$ \\
    the stream (statistics only)   & 35.71 & $-0.58$ & 91.83 & $+0.01$ \\
    structural branch              & 35.59 & $-0.70$ & 90.97 & $-0.85$ \\
    dyadic branch                  & 30.30 & $-5.99$ & 67.18 & $-24.64$ \\
    \midrule
    \emph{added} competition$^{3}$ & 36.54 & $+0.25$ & 91.62 & $-0.20$ \\
    \bottomrule
  \end{tabular}
\end{table}

\Cref{tab:ablation} removes one part at a time on ICEWS18, where the dyad-only stratum is 19\% of the queries, and on YAGO, where it is 0.1\%. Three readings follow.

\textbf{The dyadic branch is the model.} Without it \model{} is a snapshot encoder with a popularity prior and scores 30.30 on ICEWS18, which is RE-GCN's published 30.58, and 67.18 on YAGO. The branch is worth 6.0 MRR on the event dataset and 24.6 on the yearly one.

\textbf{The stream acts where the diagnostic said it would.} On ICEWS18 the Transformer over the event stream is worth 0.58 MRR over the eight summary statistics alone, and the relation types inside it are worth 0.43: a stream with types replaced by a single placeholder loses most of what the stream adds. The path intensity adds 0.37, and \cref{tab:strata} locates it: on cold-2hop queries it lifts MRR from 6.21 to 8.52 and \hten{} from 13.4 to 19.0 on seed 1. The seed spread on ICEWS18 is 0.04, so each of these effects is an order of magnitude above noise. On YAGO the same three parts are worth $+0.01$, $-0.03$ and $-0.12$, all inside the 0.20 seed spread. The eight statistics carry the dyadic branch on their own there, because on yearly data the question is only whether a fact is still running, and the count and recency of the triple answer it. This is the prediction made in \cref{sec:diag} before any model was trained, and it is the kind of evidence a stratum-level analysis can give and a leaderboard cannot.

\textbf{Two parts are null and are reported as such.} The popularity field and the candidate-context input move neither dataset outside the seed spread; a prototype intensity that scores every entity by similarity to the query's past answers, tried for the cold-far stratum, is likewise null ($+0.05$ on ICEWS18, $-0.16 \pm 0.19$ on YAGO) and is kept only as an ablation switch. We leave them in the released code with their switches rather than silently removing them, so that the ablation table is reproducible as printed.

\subsection{The same strata on every event dataset}

\begin{figure}[t]
  \centering
  % MRR of the final model by stratum on the three event datasets
% (three seeds, model without competition; from collect.py).
\begin{tikzpicture}
  \begin{axis}[
    ybar, width=\linewidth, height=5.2cm,
    bar width=5.5pt, ymin=0, ymax=100,
    ylabel={MRR ($\times 100$)}, ylabel style={font=\scriptsize},
    symbolic x coords={cold-far,cold-2hop,dyad-only,blocked,clean},
    xtick=data, tick label style={font=\scriptsize},
    enlarge x limits=0.14,
    legend style={at={(0.03,0.97)}, anchor=north west, legend columns=1,
                  font=\scriptsize, draw=none, fill=none},
    axis line style={cGrey}, tick style={cGrey},
    ymajorgrids, grid style={cGrey!25},
    nodes near coords, nodes near coords style={font=\tiny, /pgf/number format/fixed, /pgf/number format/precision=1, rotate=90, anchor=west},
  ]
    \addplot[fill=cDyad, draw=none] coordinates {
      (cold-far,1.84) (cold-2hop,8.93) (dyad-only,29.71) (blocked,34.99) (clean,85.64)};
    \addplot[fill=cDyad!55, draw=none] coordinates {
      (cold-far,2.01) (cold-2hop,11.57) (dyad-only,37.26) (blocked,38.43) (clean,89.93)};
    \addplot[fill=cGrey!60, draw=none] coordinates {
      (cold-far,0.65) (cold-2hop,2.05) (dyad-only,15.20) (blocked,21.03) (clean,80.64)};
    \legend{ICEWS18, ICEWS14s, GDELT}
  \end{axis}
\end{tikzpicture}

  \caption{\textbf{MRR by stratum on the three event datasets} (three seeds). The profile is the same on all three: the clean stratum is solved, the blocked and dyad-only strata are where a copy mechanism fails and where \model{} earns its margin, and the two cold strata are where the remaining headroom lies. GDELT, whose fifteen-minute snapshots make every stratum sparser, is lower everywhere but keeps the ordering.}
  \label{fig:strata-results}
\end{figure}

\Cref{fig:strata-results} repeats the stratified view on ICEWS14s and GDELT. The profile is identical across the three event datasets, which is what one expects of strata defined from the data rather than from the model: the clean stratum sits above 80 MRR everywhere, the blocked and dyad-only strata in the 15 to 40 range, and the cold strata near zero. The same figure says why GDELT is the hardest benchmark in the suite despite its small vocabulary: at fifteen-minute resolution its dyad streams are sparse, and every stratum except the clean one is a third to a half lower than on ICEWS.
